\section{Kiểm tra phân số là số nguyên [Trung bình]}
\subsection{Phân tích \& Thiết kế (M0)}
\begin{itemize}
    \item \textbf{Input:}
    File \texttt{fractions.csv} chứa nhiều dòng:
    \begin{itemize}
        \item Mỗi dòng có dạng: \texttt{num,den}
    \end{itemize}

    \item \textbf{Xử lý:}
    \begin{itemize}
        \item \textbf{Case biên:}
        \begin{itemize}
            \item \texttt{den = 0} $\rightarrow$ phân số không hợp lệ, bỏ qua
            \item \texttt{num = 0} $\rightarrow$ chuẩn hoá thành \texttt{0/1} và là số nguyên
            \item \texttt{den < 0} $\rightarrow$ đổi dấu cả tử và mẫu
            \item Phân số là số nguyên khi \texttt{den = 1}
        \end{itemize}

        \item \textbf{Module 1 (Đọc file):}
        \begin{itemize}
            \item Hàm \texttt{readFractions(path, vector<Fraction>\&)}
            \item Đọc từng dòng CSV
            \item Kiểm tra mở file (\texttt{fail} $\rightarrow$ \texttt{return false})
        \end{itemize}

        \item \textbf{Module 2 (Chuẩn hoá phân số):}
        \begin{itemize}
            \item Lọc invalid
            \item Đưa mẫu về dương
            \item Rút gọn gcd
        \end{itemize}

        \item \textbf{Module 3 (Kiểm tra số nguyên):}
        \begin{itemize}
            \item Sau chuẩn hoá:
            \begin{lstlisting}
den == 1
            \end{lstlisting}
        \end{itemize}

        \item \textbf{Module 4 (Ghi file):}
        \begin{itemize}
            \item Ghi \texttt{integers.csv}
            \item Chỉ ghi các phân số nguyên dạng:
            \texttt{num/1}
        \end{itemize}

        \item \textbf{Module 5 (Thống kê):}
        \begin{itemize}
            \item \texttt{total\_valid}
            \item \texttt{is\_integer\_count}
        \end{itemize}
    \end{itemize}

    \item \textbf{Output:}
    \begin{lstlisting}
Ghi file integers.csv
total_valid=<n> / is_integer_count=<m>
    \end{lstlisting}
\end{itemize}

\subsection{Mã nguồn C++11 (Tách module)}
\begin{lstlisting}
#include <iostream>
#include <vector>
#include <fstream>
#include <sstream>
#include <cstdlib>

using namespace std;

struct Fraction {
    long long num, den;
};

long long gcd(long long a, long long b) {
    a = llabs(a);
    b = llabs(b);

    while (b != 0) {
        long long t = a % b;
        a = b;
        b = t;
    }

    return a;
}

bool normalize(Fraction& f) {
    if (f.den == 0) return false;

    if (f.num == 0) {
        f.num = 0;
        f.den = 1;
        return true;
    }

    if (f.den < 0) {
        f.num = -f.num;
        f.den = -f.den;
    }

    long long g = gcd(f.num, f.den);
    f.num /= g;
    f.den /= g;

    return true;
}

bool readFractions(const string& path,
                   vector<Fraction>& arr) {
    ifstream fin(path);
    if (!fin) return false;

    string line;
    while (getline(fin, line)) {
        if (line.empty()) continue;

        stringstream ss(line);
        string s1, s2;

        if (getline(ss, s1, ',') &&
            getline(ss, s2)) {
            Fraction f;
            f.num = atoll(s1.c_str());
            f.den = atoll(s2.c_str());
            arr.push_back(f);
        }
    }

    return true;
}

bool writeIntegers(const string& path,
                   const vector<Fraction>& arr) {
    ofstream fout(path);
    if (!fout) return false;

    for (size_t i = 0; i < arr.size(); ++i)
        fout << arr[i].num
             << "/1\n";

    return true;
}

int main() {
    vector<Fraction> input, integers;

    if (!readFractions("fractions.csv", input)) {
        cerr << "Loi mo file input!\n";
        return 1;
    }

    int total_valid = 0;
    int is_integer_count = 0;

    for (size_t i = 0; i < input.size(); ++i) {
        Fraction f = input[i];

        if (normalize(f)) {
            total_valid++;

            if (f.den == 1) {
                integers.push_back(f);
                is_integer_count++;
            }
        }
    }

    if (!writeIntegers("integers.csv", integers)) {
        cerr << "Loi mo file output!\n";
        return 1;
    }

    cout << "total_valid=" << total_valid
         << " / is_integer_count="
         << is_integer_count;

    return 0;
}
\end{lstlisting}

\subsection{Kế hoạch kiểm thử (Test Cases)}

\textbf{Test Case 1 (Thông thường)}

\textbf{Input (fractions.csv):}
\begin{lstlisting}
4,2
3,1
5,2
\end{lstlisting}

\textbf{Expected Output (integers.csv):}
\begin{lstlisting}
2/1
3/1
\end{lstlisting}

\textbf{Console Output:}
\begin{lstlisting}
total_valid=3 / is_integer_count=2
\end{lstlisting}

\vspace{0.5cm}

\textbf{Test Case 2 (Có invalid)}

\textbf{Input (fractions.csv):}
\begin{lstlisting}
2,0
0,5
6,3
\end{lstlisting}

\textbf{Expected Output (integers.csv):}
\begin{lstlisting}
0/1
2/1
\end{lstlisting}

\textbf{Console Output:}
\begin{lstlisting}
total_valid=2 / is_integer_count=2
\end{lstlisting}

\vspace{0.5cm}

\textbf{Test Case 3 (Không có số nguyên)}

\textbf{Input (fractions.csv):}
\begin{lstlisting}
1,2
3,4
\end{lstlisting}

\textbf{Expected Output (integers.csv):}
\begin{lstlisting}

\end{lstlisting}

\textbf{Console Output:}
\begin{lstlisting}
total_valid=2 / is_integer_count=0
\end{lstlisting}